% Derivations checked against ../MPC/transferFunction4TankSystem.md and
% ../MPC/four-tank-model/problem4_4.m and problem4_6.m.
\chapter{Nonlinear Simulation -- Step Responses}
\label{chap:problem4}

\section{Simulate Step Response}

Step experiments characterize how each manipulated pump flow affects the
two measured lower-tank levels. All experiments use the nonlinear mass
balances and physical parameters of Chapter~\ref{chap:problem2}. The
initial condition is the equilibrium for
\begin{equation}
 u_s=[300,300]^T\ \mathrm{cm^3/s},\qquad
 \bar d=[20,20]^T\ \mathrm{cm^3/s},\qquad
 h_s=[28.4621,41.4578,4.0950,6.6332]^T\ \mathrm{cm}.
\end{equation}
Starting at equilibrium separates the input-step response from a tank-filling
transient. Pump $j$ is stepped in a separate experiment, while the other
pump remains at its nominal flow:
\begin{equation}
 u(t)=u_s+e_j\Delta F_j\,\mathbf{1}_{t\geq200},\qquad
 \Delta F_j=\alpha F_{j,s},\qquad \alpha\in\{0.10,0.25,0.50\}.
\end{equation}
Here $e_j$ is the $j$th unit vector. The increments are $30$, $75$, and
$150\ \mathrm{cm^3/s}$, giving final pump flows of $330$, $375$, and
$450\ \mathrm{cm^3/s}$. Each experiment ends at $1200\ \mathrm{s}$,
providing $1000\ \mathrm{s}$ after the step.

In the deterministic case, disturbances remain at $\bar d$ and measurement
noise is zero. MATLAB's \texttt{ode15s} integrates the intervals before and
after the step separately. Figures~\ref{fig:p4-det1} and
\ref{fig:p4-det2} show the resulting levels and pump inputs.
The direct paths $F_1\to h_1$ and $F_2\to h_2$ respond faster than the
indirect paths through the upper tanks. All positive steps increase both
lower-tank levels, consistent with the positive inflow routing.

\section{Simulate Step Response with Varying Noise Levels}

The stochastic experiments use the additive SDE from
Chapter~\ref{chap:problem2}:
\begin{equation}
 dx=f(x,u,\bar d,p)\,dt+\rho D\operatorname{diag}(s_3,s_4)\,dW,
 \qquad y_k=Cx_k+v_k,\quad v_k\sim\mathcal N(0,\sigma_y^2 I_2).
\end{equation}
Process disturbances enter the upper tanks directly, while sensor noise
is added to the lower-tank level observations. Low, medium, and high noise
scale both diffusion intensity and measurement standard deviation:
\begin{table}[htbp]
 \centering
 \caption{Noise settings for the stochastic step experiments.}
 \label{tab:p4-noise}
 \begin{tabular}{lrrr}
 \hline
 Setting & Scale & $s_3=s_4$ ($\mathrm{cm^3/s^{1/2}}$) & $\sigma_y$ (cm)\\
 \hline
 Low & 0.5 & 2.5 & 0.25\\
 Medium & 1 & 5 & 0.50\\
 High & 2 & 10 & 1.00\\
 \hline
 \end{tabular}
\end{table}
The state update is projected Euler--Maruyama with $\Delta t=0.1\ \mathrm{s}$,
using independent Brownian increments with covariance $\Delta t I_2$.
A single standardized process and measurement realization, generated with
seed 100, is reused across all step sizes and noise scales. This controls
random variation when comparing settings, but these runs are correlated
experiments rather than independent repetitions. Measurements are simulated
every $0.1\ \mathrm{s}$; the plots display them every $1\ \mathrm{s}$.
Figures~\ref{fig:p4-stoch1} and \ref{fig:p4-stoch2} distinguish true level
curves from measured level points. Increased noise causes greater measurement
scatter, while upper- and lower-tank dynamics smooth the process disturbances.

\section{Normalized Step Responses}

For input $j$ and output $i$, define the normalized response
\begin{equation}
 S_{ij}(\tau;\Delta F_j)=
 \frac{h_i(200+\tau)-h_{i,s}}{\Delta F_j},\qquad \tau\geq0.
 \label{eq:p4-normalization}
\end{equation}
For measured data, the baseline is estimated from the mean measurement over
$150\leq t<200\ \mathrm{s}$, giving
$\widehat S_{ij}=(y_i-\widehat h_{i,0})/\Delta F_j$.
This normalization gives units of $\mathrm{cm/(cm^3/s)}$ and compares the
response per unit pump increment. It does not rescale every response to
unit final value. For an LTI plant, normalized responses to different
step amplitudes would coincide.

Figures~\ref{fig:p4-normdet1}--\ref{fig:p4-normstoch2} plot normalized
responses for every pump, amplitude, and noise level. The deterministic
curves do not collapse: larger positive steps give greater normalized
steady gains. This follows directly from the equilibrium relation
$h_i=q_i^2/(2ga_i^2)$. For a lower-tank outlet flow increment
$\Delta q_i=\beta_{ij}\Delta F_j$, with
\begin{equation}
 [\beta_{ij}]=\begin{bmatrix}0.6&0.3\\0.4&0.7\end{bmatrix},
\end{equation}
the exact finite-step equilibrium gain is
\begin{equation}
 \frac{\Delta h_i^*}{\Delta F_j}
 =\frac{q_{i,s}\beta_{ij}}{ga_i^2}
  +\frac{\beta_{ij}^2\Delta F_j}{2ga_i^2}.
 \label{eq:p4-secant}
\end{equation}
The first term is the local linearized gain. The second term explains the
amplitude dependence of positive step tests.

\begin{table}[htbp]
 \centering
 \caption{Deterministic normalized levels at $1200\ \mathrm{s}$ and local
 linearized gains. All entries are in $\mathrm{cm/(cm^3/s)}$. The final
 levels approximate equilibrium; a small settling residual remains.}
 \label{tab:p4-gains}
 \begin{tabular}{lrrrr}
 \hline
 Channel & Local gain & 10\% step & 25\% step & 50\% step\\
 \hline
 $F_1\to h_1$ & 0.11777 & 0.1214 & 0.1269 & 0.1360\\
 $F_2\to h_1$ & 0.05889 & 0.0598 & 0.0612 & 0.0635\\
 $F_1\to h_2$ & 0.09476 & 0.0964 & 0.0988 & 0.1029\\
 $F_2\to h_2$ & 0.16583 & 0.1708 & 0.1782 & 0.1907\\
 \hline
 \end{tabular}
\end{table}
For fixed sensor noise, division by a larger step reduces normalized
measurement scatter: $\operatorname{SD}(v_i/\Delta F_j)=\sigma_y/\Delta F_j$.
At medium noise this is $0.0167$, $0.00667$, and $0.00333$
$\mathrm{cm/(cm^3/s)}$ for 10\%, 25\%, and 50\% steps, respectively.
Thus increasing excitation improves the signal relative to measurement
noise while increasing the departure from a local linear model.

\section{Transfer Function Identification}

The channel structure follows the analytical derivation in
Chapter~\ref{chap:problem3}: direct channels contain one tank lag and
cross channels contain two. The fitted models therefore have no zeros
or explicit delay and take the forms
\begin{equation}
 \widehat G_{ii}(s)=\frac{K_{ii}}{1+\tau_{ii}s},\qquad
 \widehat G_{ij}(s)=\frac{K_{ij}}{(1+\tau_{ij,1}s)(1+\tau_{ij,2}s)},\quad i\ne j.
\end{equation}
Identification uses separate 25\% pump steps, with medium process and
sensor noise, and seeds 101 and 102 for the two pump experiments.
Each measured baseline is estimated over the last 50 seconds before the
step. The normalized records are decimated from $0.1$ to $1\ \mathrm{s}$
sampling by selecting every tenth observation. This reduces data volume;
it is not averaging or an anti-aliasing filter.

The reported run uses the continuous-time least-squares backend.
For elapsed time $\tau\geq0$, its analytical unit-step models are
\begin{align}
 s_1(\tau)&=K(1-e^{-\tau/T}),\\
 s_2(\tau)&=K\left[1-
 \frac{T_1e^{-\tau/T_1}-T_2e^{-\tau/T_2}}{T_1-T_2}\right].
\end{align}
When $T_1=T_2=T$, the second expression has limit
$K[1-(1+\tau/T)e^{-\tau/T}]$.
MATLAB's \texttt{fminsearch} minimizes the sum of squared differences
between measured normalized post-step samples and these analytical
responses, with time constants bounded below by $0.01\ \mathrm{s}$ in
the objective evaluation. The code can alternatively use \texttt{tfest}
when System Identification Toolbox is available; its numerical estimates
may differ from those reported here.

The resulting continuous-time estimate is
\begin{equation}
 \widehat G(s)=\begin{bmatrix}
 \dfrac{0.126480}{1+83.9148s}&
 \dfrac{0.0600034}{(1+37.3589s)(1+76.0941s)}\\[6pt]
 \dfrac{0.0990653}{(1+32.6327s)(1+109.7353s)}&
 \dfrac{0.179349}{1+101.1440s}
 \end{bmatrix}.
 \label{eq:p4-identified}
\end{equation}
Rows correspond to $h_1,h_2$ and columns to $F_1,F_2$; time constants
are in seconds and gains in $\mathrm{cm/(cm^3/s)}$.
All fitted poles are real and negative. Figure~\ref{fig:p4-identification}
compares model predictions with true nonlinear normalized levels and
noisy normalized observations from the identification experiments.

\section{Accuracy of the Identified Transfer Function}

The gains in \eqref{eq:p4-identified} are close to the deterministic
25\% finite-step gains in Table~\ref{tab:p4-gains}. They represent an
approximation over that excursion, rather than the infinitesimal gain
at the initial equilibrium. The initial tank time constants are
$[74.6163,90.0541,28.3027,36.0216]^T\ \mathrm{s}$.
These explain the physical channel orders, but fitted time constants can
differ because outlet sensitivities change during a finite step and
process noise affects the observed response. The two cross-channel pole
time constants are unordered fit parameters; a fit does not uniquely
assign each pole to a particular physical tank.

To separate fit quality from generalization, the same model is evaluated
against the separate deterministic 10\%, 25\%, and 50\% experiments.
For each channel, predictions and true normalized levels are compared
on a uniform $1\ \mathrm{s}$ grid over $0\leq\tau\leq1000\ \mathrm{s}$:
\begin{equation}
 \operatorname{RMSE}_{ij}=
 \sqrt{\frac{1}{1001}\sum_{k=0}^{1000}
  \left(S_{ij}(k)-\widehat S_{ij}(k)\right)^2}.
\end{equation}
The adaptively sampled deterministic response is interpolated onto this
grid. Table~\ref{tab:p4-validation} reports the errors.
\begin{table}[htbp]
 \centering
 \caption{Normalized response RMSE against separate deterministic tests,
 in $\mathrm{cm/(cm^3/s)}$. The model is fitted to noisy 25\% experiments.}
 \label{tab:p4-validation}
 \begin{tabular}{lrrr}
 \hline
 Channel & 10\% validation & 25\% validation & 50\% validation\\
 \hline
 $F_1\to h_1$ & 0.0044 & 0.0005 & 0.0081\\
 $F_2\to h_1$ & 0.0009 & 0.0013 & 0.0029\\
 $F_1\to h_2$ & 0.0020 & 0.0007 & 0.0032\\
 $F_2\to h_2$ & 0.0072 & 0.0010 & 0.0091\\
 \hline
 \end{tabular}
\end{table}
For context, comparison against the true trajectories of the identification
runs gives normalized RMSE between $0.0010$ and $0.0012$. That is a
comparison on the fitting experiments, not independent validation.
The larger deterministic validation errors at 50\% show the limits of
using one linear model across different amplitudes. At 25\%, multiplying
normalized RMSE by $75\ \mathrm{cm^3/s}$ gives level RMSE of approximately
$0.04$--$0.10\ \mathrm{cm}$.

A useful step experiment requires an established baseline, accurate input
amplitude and timing, sufficiently long records to observe slow modes,
and excitation large enough to distinguish the response from noise.
Separate input steps isolate columns of the transfer matrix. Sampling
must resolve the fastest relevant dynamics, and model order should reflect
physical knowledge without adding unsupported parameters. Repeated runs
or separate noisy validation records would be needed to assess uncertainty
and robustness to new disturbance realizations. These results demonstrate
finite-step fit accuracy, without providing parameter confidence intervals.

\section{Markov Parameters}

The identified continuous-time model is discretized with a zero-order
hold at $T_s=4\ \mathrm{s}$, appropriate for pump commands held between
samples. For a continuous realization $(A_c,B_c,C_c,D_c)$,
\begin{equation}
 A_d=e^{A_cT_s},\qquad
 B_d=\int_0^{T_s}e^{A_c\tau}B_c\,d\tau,\qquad C_d=C_c,\quad D_d=D_c.
\end{equation}
The deviation model is
$\delta x_{k+1}=A_d\delta x_k+B_d\delta u_k$ and
$\delta z_k=C_d\delta x_k+D_d\delta u_k$.
Its Markov coefficient matrices are
\begin{equation}
 H_0=D_d,\qquad H_k=C_dA_d^{k-1}B_d\quad(k\geq1),\qquad
 \delta z_k=\sum_{\ell=0}^{k}H_\ell\delta u_{k-\ell}
 \label{eq:p4-markov}
\end{equation}
for zero initial deviation. Each $H_k$ is a $2\times2$ matrix with output
rows and input columns. The strictly proper fitted models have $H_0=0$.
For a unit input step, $S_k=\sum_{\ell=0}^kH_\ell$, so
$H_k=S_k-S_{k-1}$ for $k\geq1$. Direct differencing of noisy measured
steps would amplify noise; here the coefficients are computed from the
fitted state-space model.

Figure~\ref{fig:p4-markov} shows 60 matrices, $H_0$ through $H_{59}$,
corresponding to $0$--$236\ \mathrm{s}$. Direct-channel coefficients
are largest early and decay, while cross-channel coefficients build up
before decaying because of the upper-tank lag. These are discrete
coefficients for a unit input applied for one sample. MATLAB's
unit-area impulse convention scales the discrete response by $1/T_s$;
that scaling is not used in \eqref{eq:p4-markov}.

For stable models, the infinite coefficient sum equals the static gain:
\begin{equation}
 \sum_{k=0}^{\infty}H_k=\widehat G(0).
\end{equation}
The plotted finite horizon omits approximately
$\begin{bmatrix}6.0&8.7\\16.5&9.7\end{bmatrix}\%$
of the respective channel gains. Thus 60 coefficients illustrate the
impulse-response shapes but do not capture the entire slow response.
A longer horizon would be needed if these coefficients were used to
approximate settled outputs in a finite convolution model.

% All figures are committed report snapshots; MATLAB is not needed to build.
\begingroup
\newcommand{\pFourFigure}[1]{%
 \includegraphics[width=\linewidth]{tex/chapters/results/problem04/#1.pdf}}
\begin{figure}[htbp]
 \centering\pFourFigure{deterministic_steps_01}
 \caption{Deterministic lower-tank levels and input profiles for separate
 10\%, 25\%, and 50\% positive steps in pump 1.}
 \label{fig:p4-det1}
\end{figure}
\begin{figure}[htbp]
 \centering\pFourFigure{deterministic_steps_02}
 \caption{Deterministic lower-tank levels and input profiles for separate
 10\%, 25\%, and 50\% positive steps in pump 2.}
 \label{fig:p4-det2}
\end{figure}
\begin{figure}[htbp]
 \centering\pFourFigure{stochastic_steps_01}
 \caption{Pump 1 stochastic step tests at low, medium, and high noise.
 Curves show true levels; points show noisy observations at 1-second spacing.}
 \label{fig:p4-stoch1}
\end{figure}
\begin{figure}[htbp]
 \centering\pFourFigure{stochastic_steps_02}
 \caption{Pump 2 stochastic step tests at low, medium, and high noise.
 Curves show true levels; points show noisy observations at 1-second spacing.}
 \label{fig:p4-stoch2}
\end{figure}
\begin{figure}[htbp]
 \centering\pFourFigure{normalized_steps_01}
 \caption{Normalized deterministic responses to pump 1 steps. Differences
 between amplitudes expose the nonlinear finite-step gain.}
 \label{fig:p4-normdet1}
\end{figure}
\begin{figure}[htbp]
 \centering\pFourFigure{normalized_steps_02}
 \caption{Normalized deterministic responses to pump 2 steps.}
 \label{fig:p4-normdet2}
\end{figure}
\begin{figure}[htbp]
 \centering\pFourFigure{normalized_steps_03}
 \caption{Normalized pump 1 responses at all step sizes and noise levels.
 True curves use the equilibrium baseline; measured points use the estimated
 pre-step baseline and are displayed at 1-second spacing.}
 \label{fig:p4-normstoch1}
\end{figure}
\begin{figure}[htbp]
 \centering\pFourFigure{normalized_steps_04}
 \caption{Normalized pump 2 responses at all step sizes and noise levels,
 with true curves and noisy measured points.}
 \label{fig:p4-normstoch2}
\end{figure}
\begin{figure}[htbp]
 \centering\pFourFigure{identification_01}
 \caption{Identified channel models compared with true nonlinear and measured
 normalized responses from the noisy 25\% identification experiments.}
 \label{fig:p4-identification}
\end{figure}
\begin{figure}[htbp]
 \centering\pFourFigure{markov_01}
 \caption{Markov parameters of the fitted model under a 4-second zero-order
 hold. Coefficients beyond 236 seconds are not shown.}
 \label{fig:p4-markov}
\end{figure}
\endgroup
